\documentclass[10pt,a4paper]{amsart}
\usepackage[margin=19mm]{geometry}
\usepackage{amsmath,amssymb,mathtools}
\usepackage[scr=rsfs,cal=euler]{mathalfa}
\usepackage{xfrac}
\usepackage[unicode,hidelinks,bookmarks=false]{hyperref}
\newcommand{\ie}{i.e.\ }
\newcommand{\cf}{cf.\ }
\newcommand{\resp}{respectively\ }
\newcommand{\Subsection}{\subsection}
\providecommand{\nobreakdash}{\nobreak\hbox{-}}
\setlength{\emergencystretch}{2em}
\multlinegap0pt
\usepackage{mathtools}
\newtheorem{lemma}{Lemma}
\newtheorem{remark}{Remark}
\newtheorem{proposition}{Proposition}
\usepackage{cleveref}
\crefname{lemma}{Lemma}{Lemmas}
\crefname{remark}{Remark}{Remarks}
\crefname{proposition}{Proposition}{Propositions}
\title{Linkage axioms for generic tropical oriented matroids}
\author{Yuan Yao}
\date{Source: arXiv:2506.11265v2 (CC BY 4.0)}
\begin{document}
\maketitle

\noindent\emph{Source and licence.} Linkage axioms for generic tropical oriented matroids. Version arXiv:2506.11265v2 (CC BY 4.0), distributed by arXiv under the Creative Commons Attribution 4.0 International licence, http://creativecommons.org/licenses/by/4.0/, as recorded in the arXiv licence metadata for this exact version and shown on the abstract page https://arxiv.org/abs/2506.11265v2. Full licence notice: \texttt{licenses/arxiv-2506.11265-CC-BY-4.0.txt}.\par\smallskip

\noindent\textit{Editorial scope.} Counted unit: the article's proposition prop:treelinkagetree (source0.tex lines 679--681). The article's complete original proof of it is delivered with it (source0.tex lines 683--692, one \texttt{proof} environment of 10 lines). No mathematics is written, repaired or re-derived: the statement and the proof are the article's own lines, in the article's own notation, and no retained line is substituted or re-flowed.  Retained uncounted as the source-local context the proof needs: lines 91--93; lines 675--677.  Nothing was omitted from any retained span, so the package carries no omission record.  No retained line contains a citation, so the wrapper carries no bibliography and no bibliography entry.  No retained line contains a figure environment, an image-inclusion command or drawing code, so no image is delivered and the package declares no asset dependency.  Editorial formatting only, in addition: an AMS \texttt{amsart} preamble that restates the article's own declarations and macros used by the retained lines, namely the environments \texttt{lemma}, \texttt{remark}, \texttt{proposition} on separate counters, together with the standard packages those lines need.  The retained environments print in this document as Lemma 1, Remark 1, Proposition 1 in order of appearance, whereas the article prints the same items with the numbering of its own full document.  The complete native source of the article as distributed in the arXiv e-print for this exact version is bundled unchanged as \texttt{sources/179773/source0.tex}.
\begin{lemma}\label{lem:adjtrees}
    Suppose there are two distinct spanning trees $G, G'$ of $K_{n,d}$ and two edges $e\in G$ and $e'\in G'$ such that $G\setminus e = G' \setminus e' = H$. Let the two connected components of $H$ be $H^{(1)} = I^{(1)}\sqcup \bar{J}^{(1)}$ and $H^{(2)} = I^{(2)}\sqcup \bar{J}^{(2)}$. The two trees $G$ and $G'$ are compatible if and only if either $e$ connects $I^{(1)}$ with $\bar{J}^{(2)}$ and $e'$ connects $I^{(2)}$ with $\bar{J}^{(1)}$, or vice versa. In particular, this also means that $e$ and $e'$ do not share vertices.
\end{lemma}

\begin{remark}\label{rem:basicproperties}
    Note that the tree-linkage axiom translates to the property that for each non-leaf neighbor $i\in [n]$ of $\bar{j}$ in $\mathbb{T}(t+e_{\bar{j}})$, there is exactly one edge in $\mathbb{G}(t)$ incident to $\bar{j}$ for which the near label is $i$. (Note that since $RD(\mathbb{T}(t+e_{\bar{j}})) = t + e_{\bar{j}} + \mathbf{1}_{[\bar{d}]}$, $\bar{j}$ cannot be a leaf.) From \cref{lem:adjtrees}, the two endpoints of an edge cannot have the same label.
\end{remark}

\begin{proposition}\label{prop:treelinkagetree}
    For each lattice point $t\in (n-2)\Delta^{d-1}$, $G(t)$ is a spanning tree of $[\bar{d}]$.
\end{proposition}

\begin{proof}
    We first show that the graph is acyclic. Suppose that there is a cycle involving the vertex $\bar{j}_0$ like $\dots \bar{j}_\ell \xleftrightarrow{i_{\ell}\quad i'_0} \bar{j}_0 \xleftrightarrow{i_0\quad i'_{1}} \bar{j}_1 \dots$ in $\mathbb{G}(t)$. Denote $\mathbb{T}_k = \mathbb{T}(t+e_{\bar{j}_k})$ for convenience, then we can obtain $\mathbb{T}_1, \mathbb{T}_2, \dots, \mathbb{T}_\ell, \mathbb{T}_0 $ from $\mathbb{T}_0$ by repeatedly replacing one edge at a time. Note that the set of neighbors of $\bar{j}$ in this tree $\mathbb{T}_k$ is only modified at the first and last step, where in the first step we removed $i_0$ and in the last step we added $i_0'$. From \cref{rem:basicproperties} we see that $i_0 \neq i_0'$, which means that the final tree cannot be the same as the $\mathbb{T}_0$. Therefore, this cycle cannot exist.

    We now show that the graph is connected by showing that there is a path from $\bar{j}_0$ to $\bar{j}^*$ for any $\bar{j}_0, \bar{j}^*\in [\bar{d}]$. Suppose that in the tree $\mathbb{T}(t+e_{\bar{j}_0})$, the first vertex along the (unique) path from $\bar{j}_0$ to $\bar{j}^*$ is some $i_0\in [n]$. (In other words, $(i_0, \bar{j}_0)$ is an edge of $(\mathbb{T}(t+e_{\bar{j}_0}))^{\leftarrow\bar{j}^*}$.) This means that there is an edge of the form $\bar{j}_0 \xleftrightarrow{i_0\quad i'_1}\bar{j}_1$ in $\mathbb{G}(t)$. If $\bar{j}_1 = \bar{j}^*$ then we are done immediately. 
    Otherwise, note that from \cref{lem:adjtrees}, if $H_0 = \mathbb{T}(t+e_{\bar{j}_0})\setminus (i_0, \bar{j}_0)$, then $\bar{j}_0$ is in the same connected component $H_0^{(1)}$ of $H$ as $i'_1$, and $i_0, \bar{j}_1, \bar{j}^*$ are in the opposite connected component $H_0^{(2)}$. 
    
    Therefore, after adding the edge $(i'_1, \bar{j}_1)$, $i'_1$ is further away from $\bar{j}^*$ than $\bar{j}_1$ in the new tree $\mathbb{T}(t+e_{\bar{j}_1})$, and hence there exists a different $i_1$ adjacent to $\bar{j}_1$ that is closer to $\bar{j}^*$. If we define $H_1^{(1)}$ and $H_1^{(2)}$ as the two connected components of $H_1 = \mathbb{T}(t+e_{\bar{j}_1})\setminus (i_1, \bar{j}_1)$ where the former component contains $\bar{j}_1$, then the vertex set of $H_0^{(1)}$ is a strict subset of that of $H_1^{(1)}$, since the previous edge replacement did not affect connectivity of $H_0^{(1)}$, and this component is now also connected to $\bar{j}_1$ (and possibly some additional vertices). We also have that $\bar{j}^*\in H_1^{(2)}$ by the same argument as with $H_0$.
    
    Hence we can always repeat this process (which gives a path along the graph $\mathbb{G}(t)$), never repeating any $\bar{j}_k$'s (as $\mathbb{G}(t)$ is acyclic), until we reach $\bar{j}_\ell = \bar{j}^*$ for some $\ell$.
\end{proof}

\end{document}
